\documentclass[11pt]{article}
\usepackage[margin=1in]{geometry}
\usepackage{amsmath,amssymb,amsthm}
\usepackage{graphicx}
\usepackage{hyperref}

\newtheorem{theorem}{Theorem}
\newtheorem{lemma}{Lemma}
\newtheorem{corollary}{Corollary}
\newtheorem{definition}{Definition}
\newtheorem{remark}{Remark}

\newcommand{\ol}{\overline}
\newcommand{\Ball}{B}
\newcommand{\range}{\operatorname{range}}

\title{A shrinking-skeleton positive criterion for fully Mackey complete duals}
\author{Run \texttt{fa\_banach\_001}, agent \texttt{agent\_lane\_05}}
\date{June 14, 2026}

\begin{document}
\maketitle

\section*{Status}

This packet records a claimed partial positive answer to Problem 4.15 of
Guirao, Martinez-Cervantes and Rodriguez~\cite{GMR}.  It is not a full solution
of the problem.  It proves the desired implication for Banach spaces whose
bidual is controlled by a shrinking separable projectional skeleton.

\section*{Source problem}

Guirao, Martinez-Cervantes and Rodriguez~\cite[Problem 4.15]{GMR} ask:

\begin{quote}
Is $X^*$ fully Mackey complete whenever $X$ is $w^*$-sequentially dense in
$X^{**}$?
\end{quote}

The relevant source characterization is~\cite[Theorem 3.5]{GMR}: a Banach
space $W$ is fully Mackey complete if and only if it satisfies property $(LK)$:
for every norming norm-closed subspace $Y\subset W^*$ and every
$u\in W^*\setminus Y$, there is an absolutely convex weak-star compact
$K\subset W^*$ such that
\[
K\subset Y\oplus [u],
\qquad
u\in \ol{K\cap Y}^{\,w^*}.
\]

\begin{center}
\includegraphics[width=0.95\linewidth]{figures/open_problem_crop.png}
\end{center}

\section*{Proof Intuition}

The previous $\ell_p$-sum packet showed that countable coordinate supports let
one localize a norming subspace $Y\subset X^{**}$ to a separable countable
sub-sum.  The present packet isolates the abstract mechanism behind that
argument.

Suppose $X$ has many separable projection ranges $P_sX$, every point of
$X^{**}$ is fixed by some $P_s^{**}$, and the projections are shrinking along
countable suprema.  Given a norming closed $Y\subset X^{**}$ and a starting
range containing $u$, recursively do the following.  On the current separable
range, the dual ball is weak-star compact metrizable.  A countable subcover
therefore selects countably many elements of $Y$ which norm that current range.
Since every selected element of $Y$ is captured somewhere by the skeleton,
move to a later range containing all of them.  At the countable supremum, the
selected elements lie in $Y$ intersected with the final bidual range, and the
shrinking hypothesis passes the norming estimates to the whole final dual.

The source paper's separable theorem then applies to the final separable range:
it contains no $\ell_1$, so its dual is fully Mackey complete.  Its $LK$
witness embeds back into $X^{**}$ because the final range is the range of a
projection.

\section*{Definitions}

Let $(\Sigma,\leq)$ be a directed partially ordered set in which every
increasing sequence has a supremum.

\begin{definition}
A family $(P_s)_{s\in\Sigma}$ of bounded projections on a Banach space $X$ is
called a shrinking bidual-capturing separable projectional skeleton if:
\begin{enumerate}
\item $\sup_s\|P_s\|<\infty$ and $P_sX$ is separable for every $s$;
\item if $s\leq t$, then $P_sP_t=P_tP_s=P_s$;
\item if $s=\sup_n s_n$ for an increasing sequence $(s_n)$, then
$P_sX=\ol{\bigcup_n P_{s_n}X}$ and, for every $\phi\in (P_sX)^*$,
\[
\|\phi-\phi\circ P_{s_n}|_{P_sX}\|_{(P_sX)^*}\longrightarrow 0;
\]
\item for every $x^{**}\in X^{**}$ there is $s\in\Sigma$ such that
$P_s^{**}x^{**}=x^{**}$.
\end{enumerate}
\end{definition}

\section*{Main theorem}

\begin{theorem}\label{thm:main}
Let $X$ be a Banach space containing no subspace isomorphic to $\ell_1$.
Assume that $X$ admits a shrinking bidual-capturing separable projectional
skeleton.  Then $X$ is weak-star sequentially dense in $X^{**}$, and $X^*$ is
fully Mackey complete.
\end{theorem}

\section*{Localization lemma}

For $s\in\Sigma$, put $E_s=P_sX$ and $R_s=P_s^{**}X^{**}$.  We identify $R_s$
with $E_s^{**}$ in the usual way.  The weak-star topology induced from $X^{**}$
on $R_s$ agrees with $\sigma(E_s^{**},E_s^*)$, since every member of $E_s^*$
extends to $X^*$.

\begin{lemma}\label{lem:localization}
Let $Y\subset X^{**}$ be a norming norm-closed subspace.  For every
$s_0\in\Sigma$ there is $s\geq s_0$ such that
\[
Y_s:=Y\cap R_s
\]
is norming for $E_s^*$.
\end{lemma}

\begin{proof}
Let $M=\sup_s\|P_s\|$.
Choose $0<\delta\leq1$ such that
\[
\sup\{|y(x^*)|:\ y\in Y,\ \|y\|\leq1\}\geq \delta\|x^*\|
\qquad (x^*\in X^*).
\tag{1}
\]

We recursively construct an increasing sequence $s_0\leq s_1\leq s_2\leq
\cdots$ and countable sets $C_n\subset Y\cap\Ball_{X^{**}}$.  Suppose $s_n$
has been chosen.  Since $E_{s_n}$ is separable, the compact space
$(\Ball_{E_{s_n}^*},w^*)$ is metrizable.  For each rational $r\in(0,1)$, set
\[
O_r=\{\phi\in\Ball_{E_{s_n}^*}:\ \|\phi\|>r\}.
\]
For $\phi\in O_r$, apply (1) to $P_{s_n}^*\phi=\phi\circ P_{s_n}\in X^*$.
Because $P_{s_n}^*\phi$ restricts to $\phi$ on $E_{s_n}$, we have
$\|P_{s_n}^*\phi\|\geq\|\phi\|>r$.  Hence some
$y\in Y\cap\Ball_{X^{**}}$ satisfies
\[
|P_{s_n}^{**}y(\phi)|=|y(P_{s_n}^*\phi)|>\delta r.
\]
Thus the weak-star open sets
\[
\{\psi\in\Ball_{E_{s_n}^*}: |P_{s_n}^{**}y(\psi)|>\delta r/2\},
\qquad y\in Y\cap\Ball_{X^{**}},
\]
cover $O_r$.  Since $O_r$ is Lindelof, choose a countable subcover.  Taking
the union over rational $r\in(0,1)$ gives a countable set
$C_n\subset Y\cap\Ball_{X^{**}}$ such that
\[
\sup_{y\in C_n}|P_{s_n}^{**}y(\phi)|
\geq \frac{\delta}{2}\|\phi\|
\qquad(\phi\in E_{s_n}^*).
\tag{2}
\]

For every $y\in C_n$, choose $t_y\in\Sigma$ such that $P_{t_y}^{**}y=y$.
Using directedness and countable suprema, choose $s_{n+1}\geq s_n$ which
dominates all $t_y$ with $y\in C_n$.  Then $C_n\subset Y\cap R_{s_{n+1}}$.

Let $s=\sup_n s_n$.  Then $C_n\subset Y\cap R_s$ for every $n$.  We claim that
$Y_s$ is norming for $E_s^*$.  Let $0\neq\phi\in E_s^*$, and write
\[
\phi_n:=\phi\circ P_{s_n}|_{E_s}.
\]
By the shrinking assumption, $\|\phi-\phi_n\|\to0$.  Also
$\phi_n|_{E_{s_n}}\in E_{s_n}^*$, and
\[
\|\phi_n\|\leq M\|\phi_n|_{E_{s_n}}\|.
\]
Fix a small $\eta>0$ and choose $n$ so large that
\[
\|\phi-\phi_n\|<\eta\|\phi\|,
\qquad
\|\phi_n|_{E_{s_n}}\|>\frac{1-\eta}{M}\|\phi\|.
\]
By (2), choose $y\in C_n$ within $\eta\|\phi\|$ of the supremum at
$\phi_n|_{E_{s_n}}$.  Since $y\in R_s$ and $\|y\|\leq1$,
\[
\begin{aligned}
|y(\phi)|
&\geq |y(\phi_n)|-|y(\phi-\phi_n)|\\
&= |P_{s_n}^{**}y(\phi_n|_{E_{s_n}})|-|y(\phi-\phi_n)|\\
&> \left(\frac{\delta}{2M}(1-\eta)-2\eta\right)\|\phi\|.
\end{aligned}
\]
Taking $\eta$ sufficiently small gives
\[
\sup\{|y(\phi)|:\ y\in Y_s,\ \|y\|\leq1\}\geq \frac{\delta}{4M}\|\phi\|.
\]
Therefore $Y_s$ is norming for $E_s^*$.
\end{proof}

\section*{Proof of Theorem \ref{thm:main}}

First let $u\in X^{**}$.  Choose $s$ with $P_s^{**}u=u$.  The range $E_s$ is
separable and contains no $\ell_1$, because it is a subspace of $X$.  By
\cite[Corollary 4.8]{GMR}, $(\Ball_{E_s^{**}},w^*)$ is Frechet-Urysohn.  Since
$E_s$ is weak-star dense in $E_s^{**}$ by Goldstine's theorem, there is a
sequence in $E_s\subset X$ converging weak-star to $u$.  Hence $X$ is
weak-star sequentially dense in $X^{**}$.

It remains to show that $X^*$ is fully Mackey complete.  We verify property
$(LK)$ from~\cite[Theorem 3.5]{GMR}.  Let $Y\subset X^{**}$ be norming and
norm-closed, and let $u\in X^{**}\setminus Y$.  Choose $s_0$ with
$P_{s_0}^{**}u=u$.  By Lemma~\ref{lem:localization}, after increasing $s_0$ if
necessary, we may assume that $Y_s=Y\cap R_s$ is norming for $E_s^*$ and still
$u\in R_s$.

The space $E_s$ is separable and contains no $\ell_1$, so
\cite[Corollary 4.8]{GMR} gives that $E_s^*$ is fully Mackey complete.
Applying property $(LK)$ in the dual space $E_s^*$ to the norming closed
subspace $Y_s\subset E_s^{**}$ and to $u\notin Y_s$, we obtain an absolutely
convex weak-star compact set $K\subset R_s$ such that
\[
K\subset Y_s\oplus [u],
\qquad
u\in \ol{K\cap Y_s}^{\,\sigma(E_s^{**},E_s^*)}.
\]
Regarded as a subset of $X^{**}$, the same $K$ is weak-star compact, because
the weak-star topology on $R_s$ is the one induced from $X^{**}$.  Moreover
$Y_s\subset Y$, and therefore
\[
K\subset Y\oplus [u],
\qquad
u\in \ol{K\cap Y}^{\,\sigma(X^{**},X^*)}.
\]
This is the required $LK$ witness.  Thus $X^*$ is fully Mackey complete.

\section*{Relation to the earlier \texorpdfstring{$\ell_p$}{ell-p}-sum packet}

The criterion recovers the previous positive theorem for
\[
X=\left(\bigoplus_{\gamma\in\Gamma}E_\gamma\right)_{\ell_p},
\qquad 1<p<\infty,
\]
when each $E_\gamma$ is separable and contains no $\ell_1$.  Take $\Sigma$ to
be the countable subsets of $\Gamma$ and let $P_A$ be the coordinate
projection.  Every element of $X^{**}=(\oplus_\gamma E_\gamma^{**})_{\ell_p}$
has countable support, so it is captured by some $P_A^{**}$.  Along increasing
countable sets, the dual tails in
$X^*=(\oplus_\gamma E_\gamma^*)_{\ell_q}$ converge in norm.  Finally, $X$
contains no $\ell_1$: any bounded sequence is supported in a countable sub-sum,
and the standard Rosenthal diagonal argument applies there.

\section*{Verification and limitations}

This is a partial positive result, not a full solution of Problem 4.15.  It
identifies a sufficient structural condition under which the countable
localization argument works.  The proof does not show that every Banach space
which is weak-star sequentially dense in its bidual has such a skeleton.  In
fact, that is exactly the obstruction exposed by the argument: weak-star
sequential density supplies separable weak-star closures of individual
bidual points, but it does not by itself supply projections whose adjoints
converge in norm on the relevant separable duals.

The main point for review is Lemma~\ref{lem:localization}, especially the last
passage from the estimates on $E_{s_n}^*$ to the estimate on $E_s^*$.  This is
where the shrinking condition is used and where a full solution would need a
replacement mechanism.

\begin{thebibliography}{9}

\bibitem{GMR}
A. J. Guirao, G. Martinez-Cervantes, and J. Rodriguez,
\emph{Completeness in the Mackey topology by norming subspaces},
J. Math. Anal. Appl. \textbf{478} (2019), no. 2, 776--789;
arXiv:1812.10669.

\end{thebibliography}

\end{document}
